% Chapitre 2 adapté du fichier Word corrigé fourni par l'auteur.
% Le tableau récapitulatif de la revue théorique et empirique est volontairement exclu.

\newcommand{\concept}[1]{%
  \par\medskip\noindent
  \professionalbullet\hspace{0.65em}%
  {\color{navyblue}\bfseries #1}%
  \par\nobreak\smallskip
}

\chapter{CADRE CONCEPTUEL ET REVUE DE LITTÉRATURE}

\begin{chapterintro}
\begin{justify}
Ce chapitre a pour objectif de préciser les principaux concepts nécessaires à la compréhension des dynamiques climatiques et de la variabilité pluviométrique. Il propose ensuite une revue des contributions théoriques et empiriques consacrées à l'évolution des précipitations, aux ruptures climatiques et aux régimes pluviométriques, afin de situer l'étude dans la littérature existante.
\end{justify}
\end{chapterintro}

\vspace{0.45cm}
\noindent{\color{navyblue}\rule{\textwidth}{1.1pt}}
\vspace{0.35cm}

\section{Cadre conceptuel}

\subsection{Climat}

Le mot « climat » vient du grec « klima », qui signifie « inclinaison ». Étymologiquement, il désigne la hauteur du Soleil en un lieu et à une saison donnée, traduisant l'ancrage originel de la climatologie dans l'observation astronomique plutôt que météorologique.

Le climat désigne également la synthèse statistique des conditions atmosphériques observées en un lieu donné sur une période de référence, généralement fixée à trente ans par l’Organisation météorologique mondiale (OMM, 2018). Il fait ainsi référence à l'état moyen de l'atmosphère, incluant la température, la pluviométrie, l'humidité et le vent. Il se distingue en cela du temps, qui renvoie à l'état instantané de l'atmosphère à un moment donné. Cet état moyen, considéré dans sa globalité, constitue le point de départ des définitions plus spécifiques présentées ci-après.

\subsection{Dynamique climatique}

La dynamique climatique désigne l'ensemble des processus physiques qui régissent le fonctionnement du climat à travers les échanges d'énergie entre l'atmosphère, l'océan et la surface continentale (GIEC, 2021). Elle se distingue ainsi de la simple description statistique du climat évoquée précédemment : alors que celui-ci renvoie à un état moyen, la dynamique climatique s'intéresse aux mécanismes qui produisent cet état et ses variations.

Dans l'espace sahélien, elle renvoie principalement au fonctionnement de la mousson ouest-africaine, dont l'établissement saisonnier commande l'essentiel du régime pluviométrique national, ainsi qu'aux mouvements de grande échelle de l'atmosphère et des océans, qui régulent les transferts de chaleur entre les basses et les hautes latitudes.

\concept{La circulation atmosphérique générale}
Elle représente le mouvement organisé des masses d'air à l'échelle planétaire sous l'effet du réchauffement inégal de la Terre entre l'équateur et les pôles. Ce déséquilibre thermique engendre des cellules de convection\footnote{Cellule de convection : C'est une circulation en boucle où l’air chaud s’élève, se refroidit en altitude puis redescend.}, dont la cellule de Hadley, qui achemine l'air chaud et humide de l'équateur vers les tropiques, où il redescend en formant les hautes pressions subtropicales caractéristiques des zones désertiques.

\concept{La zone de convergence intertropicale (ZCIT)}
Elle désigne la bande de basses pressions située près de l’équateur, où convergent les alizés des deux hémisphères. L’air chaud et humide s’y élève, ce qui produit une forte nébulosité et des précipitations abondantes. Elle se déplace avec le Soleil : plus au nord pendant l’été, plus au sud pendant l’hiver. Au Sénégal, sa remontée vers le nord de juin à septembre apporte la saison des pluies.

\concept{La mousson ouest-africaine}
Elle correspond au mécanisme saisonnier d'inversion des vents et des flux d'humidité, associé au déplacement de la zone de convergence intertropicale (ZCIT) entre les deux hémisphères. Son installation, généralement entre juin et septembre au Sénégal, commande la quasi-totalité des précipitations annuelles enregistrées sur le territoire.

\concept{Les interactions océan--atmosphère}
Elles désignent les échanges de chaleur et d'humidité entre la surface océanique et l'atmosphère, qui modulent la circulation atmosphérique tropicale. Les océans, en stockant et en restituant la chaleur avec une forte inertie, influencent la trajectoire des systèmes pluvieux bien au-delà de leurs propres limites géographiques.

\concept{Le forçage climatique}
Il représente toute perturbation, externe ou interne au système climatique, susceptible de modifier son équilibre énergétique (GIEC, 2021). Il peut être d'origine astronomique, radiative, océanique ou anthropique. Autrement dit, tout élément capable de perturber les entrées ou les sorties d'énergie du système climatique constitue un forçage climatique, qu'il s'agisse d'une modification du rayonnement solaire reçu, de la composition de l'atmosphère ou des échanges océan--atmosphère.

\subsection{Variabilité climatique}

La variabilité climatique désigne les fluctuations naturelles du climat autour de sa valeur moyenne, à des échelles de temps variées (Hulme, 1992). Elle concerne aussi bien des échelles courtes, telles que les saisons ou les années, que des échelles plus longues, de l'ordre de la décennie. Le GIEC (2021) précise que cette variabilité renvoie à la variance interne à l'intérieur d'un même régime climatique. Une phase sèche prolongée peut ainsi relever d'une simple variabilité interne du système climatique sahélien, sans traduire nécessairement une modification durable de ses paramètres moyens.

Cette variabilité trouve son origine dans trois grandes sources. D’abord, une variabilité naturelle externe, liée à des forçages extérieurs au système climatique tels que le forçage solaire, les cycles de Milankovitch ou le volcanisme. Ensuite, une variabilité naturelle interne, propre aux rétroactions internes du système climatique, notamment le couplage océan--atmosphère (ENSO, NAO, circulation thermohaline)\footnote{\textit{ENSO} : El Niño--Oscillation australe. C’est un phénomène couplé océan--atmosphère du Pacifique tropical qui alterne, de façon irrégulière, entre une phase chaude (El Niño) et une phase froide (La Niña). \textit{NAO} : Oscillation nord-atlantique. Il désigne la fluctuation de la différence de pression entre l’anticyclone des Açores et la dépression d’Islande, qui module les vents et les perturbations sur l’Atlantique Nord. \textit{Circulation thermohaline} : C’est la circulation océanique de grande échelle, engendrée par les différences de densité des eaux (température et salinité), qui transporte la chaleur entre les basses et les hautes latitudes.}. Enfin, une variabilité anthropique externe, résultant des activités humaines, à travers la perturbation de l'effet de serre et des conditions de surface.

On distingue notamment les formes suivantes.

\concept{La variabilité interannuelle}
Elle correspond à la succession d’années excédentaires et déficitaires en précipitations, d’une année à l’autre. Cette forme de variabilité est particulièrement marquée dans le régime climatique sahélien, où l’écart entre l’année la plus pluvieuse et l’année la plus sèche d’une même série peut dépasser 700mm.

\concept{La variabilité décennale}
Elle désigne l'alternance de phases pluriannuelles globalement sèches ou humides, telle que la longue période de déficit pluviométrique observée au Sahel entre la fin des années 1960 et le début des années 2000.

\subsection{Changement climatique}

Le changement climatique désigne une modification durable de l’état du climat, identifiable par des changements de sa moyenne et/ou de sa variabilité, qui persiste sur une longue période, généralement plusieurs décennies ou plus (GIEC, 2021). Il se distingue ainsi de la variabilité climatique, qui correspond aux fluctuations naturelles et à plus court terme autour de cet état moyen.

Ces modifications, qui durent généralement plusieurs décennies, se traduisent par des tendances progressives ou par des ruptures statistiques dans les séries observées. Le changement climatique peut résulter de forçages naturels, tels que les variations de l'activité solaire ou volcanique, ou de forçages anthropiques, au premier rang desquels figure l'augmentation des concentrations atmosphériques de gaz à effet de serre. L'influence de ces derniers sur le climat contemporain est aujourd'hui documentée avec un niveau de confiance élevé par les rapports successifs du GIEC.

Alors que la variabilité climatique renvoie à des fluctuations réversibles autour d'un état moyen stable, le changement climatique implique le passage à un nouvel état d'équilibre, sans garantie que le système revienne à sa configuration antérieure. Une rupture statistiquement significative ne suffit donc pas, à elle seule, à établir l'existence d'un changement climatique au sens strict du terme.

\subsection{Régime pluviométrique et indicateurs de mesure}

Le régime pluviométrique désigne les caractéristiques de la répartition saisonnière des précipitations en un lieu donné (OMM, 2012). Il détermine la durée et l’intensité des saisons sèches et humides. Au Sénégal, l’hivernage, qui s’étend généralement de juin à octobre, constitue la principale saison des pluies. Le pays présente un gradient pluviométrique national marqué, la pluviométrie moyenne annuelle augmentant du nord vers le sud (Touré et al., 2022). La pluviométrie, expression la plus concrète et la plus directement mesurable des concepts précédents, constitue la variable d’observation retenue dans le présent travail, à travers ses cumuls annuels. Par extension, on désigne ici par « régimes pluviométriques », les grands états pluviométriques interannuels (années sèches, normales et humides) qui se dégagent de l’analyse.

\concept{Les systèmes convectifs de méso-échelle (SCM)}
Ce sont de vastes ensembles d'orages organisés qui génèrent l'essentiel des précipitations en zone sahélienne. Les lignes de grains\footnote{Ligne de grains : ligne d’orages organisés en bande, accompagnée de vents et de pluies intenses au passage.}, caractérisées par un front convectif linéaire\footnote{Front convectif linéaire : Alignement d’orages formés par convection (ascension de l’air chaud et humide) et qui se déplace sous forme de bande.} pouvant s'étendre sur plusieurs centaines de kilomètres, en constituent l'exemple le plus typique en Afrique de l'Ouest.

\concept{L’indice d’aridité}
Il évalue le degré de sécheresse d’un milieu en comparant les précipitations annuelles à la température ou à l’évapotranspiration potentielle. Les indices de De Martonne et de Gaussen figurent parmi les plus couramment mobilisés dans la caractérisation climatique de l’Afrique de l’Ouest.

\paragraph{Indice de De Martonne}
L’indice de De Martonne est calculé selon la formule :

\begin{formulabox}[0.42\linewidth]{}
\[
I = \frac{P}{T+10}
\]
\end{formulabox}

Dans cette formule, $P$ désigne la précipitation totale annuelle, exprimée en mm, et $T$ la température moyenne annuelle, exprimée en °C. De Martonne a étudié la relation entre la valeur de cet indice et le type de végétation observé afin de caractériser le climat général d’une région.

\begin{table}[H]
\centering
\captionsetup{justification=centering,font=small,labelfont=bf}
\caption{Classes climatiques associées à l’indice d’aridité de De Martonne}
\label{tab:classes-indice-de-martonne}
\begingroup
\footnotesize
\renewcommand{\arraystretch}{1.18}
\setlength{\tabcolsep}{4pt}
\begin{tabularx}{\linewidth}{>{\centering\arraybackslash}p{1.7cm}>{\raggedright\arraybackslash}p{2.3cm}>{\raggedright\arraybackslash}p{3cm}X}
\toprule
\rowcolor{navyblue}
\color{white}\textbf{Valeur de $I$} &
\color{white}\textbf{Type de climat} &
\color{white}\textbf{Végétation potentielle} &
\multicolumn{1}{>{\centering\arraybackslash}X}{\color{white}\textbf{Exemples}} \\
\midrule
\rowcolor{navyblue!7}
0 à 5 & Hyperaride & Désert absolu & Vallée de la Mort (États-Unis) \\
5 à 10 & Aride & Désert & Désert du Sahara ; désert de l’Arizona (États-Unis) \\
\rowcolor{navyblue!7}
10 à 20 & Semi-aride & Steppe & Sahel, Kalahari \\
20 à 30 & Semi-humide & Prairie naturelle, forêt & Région parisienne, Orléans, région de Marseille \\
\rowcolor{navyblue!7}
30 à 40 & Humide & Forêt & Région de Lille \\
40 à 55 & Humide & Forêt & Rouen, Dieppe, Alençon \\
\bottomrule
\end{tabularx}
\endgroup
\par\smallskip
{\footnotesize\itshape Source : tableau fourni par l’auteur.}
\end{table}

\paragraph{Indice de Gaussen ou indice ombrothermique}
L’indice de Gaussen compare, mois par mois, les précipitations $P$ (en mm) à deux fois la température moyenne $T$ (en °C). D’après F. Bagnouls et H. Gaussen (1953), un mois est considéré comme sec lorsque ses précipitations sont inférieures ou égales à deux fois sa température moyenne, soit :

\begin{formulabox}[0.30\linewidth]{}
\[
P \leq 2T
\]
\end{formulabox}

Dans cette relation, $P$ désigne les précipitations du mois et $T$ sa température moyenne. Bagnouls et Gaussen ont développé le diagramme ombrothermique pour représenter cette comparaison au cours de l’année.


\concept{Le cumul pluviométrique et le nombre de jours de pluie}
Le cumul pluviométrique mesure la quantité totale de précipitations reçues sur une période donnée, qu'elle soit mensuelle, saisonnière ou annuelle. Le nombre de jours de pluie dénombre, quant à lui, les jours au cours desquels un seuil minimal de précipitation a été enregistré. La confrontation de ces deux indicateurs permet de nuancer l'interprétation d'un simple cumul : une pluviométrie stable ou en hausse, combinée à une baisse du nombre de jours pluvieux, signale une intensification des événements plutôt qu'un retour pur et simple aux conditions antérieures.

\section{Revue de littérature}

\subsection{Revue théorique}

Trois approches structurent la littérature sur les dynamiques climatiques et leur variabilité pluviométrique au Sénégal : la théorie des interactions océan--atmosphère, la théorie du bilan radiatif et de l'effet de serre, ainsi que l'approche statistique de détection des ruptures et des tendances. Ces approches ne sont pas concurrentes, mais complémentaires, chacune éclairant une dimension différente du problème tout en présentant ses propres limites.

\subsubsection{Théorie des interactions océan--atmosphère}

Cette théorie analyse les échanges de chaleur entre les océans et l'atmosphère. En raison de leur forte inertie thermique, les océans stockent la chaleur sur de longues périodes et influencent la trajectoire des vents et des systèmes pluvieux bien au-delà de leur emprise géographique propre. Appliquée au Sahel, cette approche a été formalisée par Folland, Palmer et Parker (1986), qui montrent que les phases sèches et humides observées dans la bande sahélienne sont fortement liées à des configurations contrastées des températures de surface des océans à l'échelle quasi mondiale. Ces anomalies thermiques marines provoquent des déplacements de la zone de convergence intertropicale et, par conséquent, des variations directes du régime pluviométrique régional.

Une déclinaison complémentaire, développée par Charney (1975), introduit une échelle plus locale en postulant un mécanisme de rétroaction entre la végétation, le sol et l'atmosphère. La dégradation du couvert végétal augmenterait l'albédo\footnote{Albédo : C'est une fraction du rayonnement solaire incident réfléchie par une surface, comprise entre 0 et 1. L’albédo d’un sol dépend de sa couleur, de son humidité et de son couvert végétal.} des sols, provoquerait une subsidence atmosphérique et réduirait les précipitations. Cette interprétation, longtemps mobilisée pour expliquer la sécheresse sahélienne par les pratiques agricoles locales, a depuis été nuancée par des travaux montrant que la rétroaction positive entre la végétation et les pluies relèverait davantage d'un recyclage local de l'humidité que du seul effet d'albédo initialement avancé.

\subsubsection{Théorie du bilan radiatif et de l'effet de serre}

Le climat dépend de l'équilibre entre l'énergie solaire reçue par la Terre et l'énergie qu'elle réémet vers l'espace sous forme de rayonnement infrarouge. L'effet de serre naturel, dû à la présence de gaz atmosphériques tels que la vapeur d'eau et le dioxyde de carbone, permet de maintenir une température moyenne d'environ 15°C à la surface du globe, contre une valeur théorique négative en son absence. Cette théorie, dont les fondements remontent aux travaux de Tyndall (1861), puis d'Arrhenius (1896), établit que toute augmentation de la concentration atmosphérique en gaz à effet de serre rompt l'équilibre radiatif et entraîne un réchauffement global. Elle constitue le socle théorique du changement climatique anthropique contemporain, tel que documenté par les rapports successifs du GIEC.

Ce réchauffement agit directement sur le cycle de l'eau. D'après la relation de Clausius-Clapeyron, la capacité de l'atmosphère à contenir de la vapeur d'eau augmente d'environ 7\% par degré de réchauffement. Plus d'humidité disponible alimente plus d'énergie latente dans les systèmes convectifs, d'où l'hypothèse d'une intensification du cycle hydrologique : évaporation accrue, précipitations plus intenses lors des événements pluvieux, et contrastes plus marqués entre périodes humides et périodes sèches (Held et Soden, 2006 ; Trenberth, 2011). Cette intensification n'est toutefois pas uniforme dans l'espace. En Afrique de l'Ouest, la pluie dépend aussi de la mousson, des températures de surface océanique et de la position de la zone de convergence intertropicale (Nicholson, 2013). Le signal radiatif peut donc se combiner à la variabilité régionale, ou la masquer.

Cette théorie permet de formuler l'hypothèse suivante selon laquelle une partie des changements observés dans le régime pluviométrique du Sénégal pourrait s'inscrire dans cette dynamique globale, plutôt que résulter de la seule variabilité naturelle régionale.

\subsubsection{Approche statistique de détection des ruptures et des tendances}

Cette approche ne propose pas un mécanisme causal, mais fournit un ensemble d’outils permettant de caractériser objectivement la variabilité observée dans une série chronologique. Inspirée des travaux fondateurs de Lamb (1982) sur les indices d’anomalies, puis formalisée en climatologie par Hulme (1992) et Ozer, Bodart et Tychon (2005), elle repose sur des tests de rupture (Pettitt 1979, Buishand 1982, Hubert et al. 1989, Lubès-Niel et al. 1998) ainsi que sur des tests de tendance, tel que celui de Mann–Kendall.

Ces outils permettent d’identifier les changements de régime et les évolutions tendancielles au sein des séries climatiques, tout en offrant un cadre statistique adapté à l’analyse de leur variabilité. Cette approche présente toutefois une limite fondamentale largement reconnue dans la littérature : elle demeure essentiellement descriptive, dans la mesure où elle permet d’identifier l’existence, la localisation ou l’évolution d’un changement dans une série pluviométrique sans expliquer les mécanismes physiques ou climatiques qui en sont à l’origine.

\subsection{Revue empirique}

L'étude des dynamiques climatiques au Sénégal a donné lieu à un nombre croissant de travaux empiriques qui peuvent être regroupés en trois ensembles. Le premier porte sur la rupture pluviométrique de la fin des années 1960, le deuxième documente la reprise pluviométrique partielle des années 2000 et le troisième analyse les répercussions hydrologiques et écologiques de cette variabilité.

Les études empiriques montrent de façon convergente l'existence d'une rupture pluviométrique majeure à la fin des années 1960 au Sénégal. Diop (2023) met en évidence, à l’aide de la segmentation de Hubert, les principales ruptures d’une série longue de 122 ans, couvrant la période 1901–2022 au Sénégal et en Gambie. Il établit que les sécheresses du début des années 1970 et 1980 demeurent exceptionnelles au regard de l’ensemble de la période étudiée. Il confirme également le caractère représentatif des années 1968 et 2005 dans l’évolution de la pluviométrie nationale.

Touré et al. (2022), à partir de dix stations réparties sur l'ensemble du territoire sénégalais et couvrant la période 1960-2010, mobilisent plusieurs outils statistiques complémentaires : le test séquentiel de Mann--Kendall, le test de Van Belle et Hughes et le test régional de Mann--Kendall. Ils ne détectent qu'une seule tendance significative à la baisse, à la station de Matam. Ce résultat traduit un comportement globalement irrégulier de la pluviométrie sur le reste du territoire depuis 1960.

Bodian (2014), à partir d'un échantillon de 22 stations couvrant la période 1940-2013, confirme l'existence de cette rupture en appliquant trois méthodes de détection : le test de \textsc{Pettitt}, la méthode de \textsc{Lee} et \textsc{Heghinian} et la segmentation de \textsc{Hubert}. L'auteur conclut que le déficit pluviométrique persistait encore à la fin de la période étudiée, malgré les années humides ponctuelles de 1994 et 1999. Sagna et al. (2015) situent plus précisément cette rupture entre 1966 et 1969, l'année 1968 constituant le point de bascule le plus fréquemment retenu. Gaye (2017), sur les séries de Saint-Louis, Podor et Matam couvrant la période 1954-2013, aboutit à une conclusion similaire. La rupture est située en 1968 et jusqu'à 33 des 60 années étudiées se révèlent déficitaires selon les stations considérées.

Faye (2019), dans une étude consacrée à la région de Dakar entre 1960 et 2013, à l’aide du test de \textsc{Pettitt}, du test de Mann–Kendall et des indices de tendance et de variation, ne détecte aucune tendance pluviométrique statistiquement significative sur l’ensemble de la période, malgré une diminution de 211mm, soit 36,2\%, des totaux pluviométriques observés à partir de 1969. Ce résultat illustre la différence entre une rupture ponctuelle et une tendance linéaire. La même étude relève parallèlement une rupture thermique nette en 1994, accompagnée d'une augmentation de 0,9°C de la température moyenne entre les sous-périodes 1960-1994 et 1995-2013. Le réchauffement local ne suit donc pas nécessairement le même calendrier que l'évolution des précipitations.

Depuis la fin des années 1990, la pluviométrie sahélienne montre une amélioration partielle, avec l'apparition de plusieurs années excédentaires à partir de 2005. Cette reprise, documentée notamment par Gaye (2017), ne signifie cependant pas un simple retour aux conditions antérieures à 1968. Mendy (2023), à partir de la série centennale de Dakar-Yoff, confirme ce retour pluviométrique tout en soulignant qu'il coïncide avec une recrudescence des inondations urbaines dans l'agglomération dakaroise. Cela suggère une intensification des événements pluvieux plutôt qu'une simple réversibilité du régime antérieur. Cette interprétation rejoint les résultats obtenus à l'échelle nationale par Diop (2023), qui relève une diminution du nombre de jours de pluie malgré la reprise des cumuls. La sécheresse persisterait ainsi sous une forme différente de celle des décennies 1970-1990.

Un dernier ensemble de travaux documente les répercussions hydrologiques et écologiques de cette variabilité climatique, avec des résultats qui convergent avec ceux obtenus sur la pluviométrie. Dans le delta du fleuve Sénégal, Ngom (2013), en mobilisant les mêmes méthodes de détection de rupture que Bodian (2014), montre que la sécheresse a fortement réduit les écoulements naturels du fleuve. La régulation assurée par les barrages de Diama et de Manantali a depuis partiellement sécurisé l’approvisionnement en eau douce de la basse vallée.

Plus au sud, Faye (2017) analyse les débits du bassin de la Falémé à l'aide du test de Mann--Kendall appliqué aux séries mensuelles et annuelles couvrant la période 1954-2014. Son analyse révèle une baisse statistiquement significative des écoulements, atteignant 50\% à 59\% selon les stations de Fadougou, Gourbassi et Kidira. Ce résultat corrobore, au moyen d'un indicateur indépendant, l'ampleur de la péjoration pluviométrique documentée par ailleurs.

Sur la frange littorale des Niayes, les travaux d'Aguiar (2009) mettent en évidence une diminution progressive de la nappe phréatique côtière consécutive à la baisse des précipitations, favorisant une intrusion saline responsable du dépérissement du palmier à huile et d'autres espèces sensibles à la salinité des sols. Plus récemment, Sow et al. (2025), dans le sud du pays, relient directement la variabilité pluviométrique à la dégradation du couvert végétal de la forêt classée de Pata, à partir d'une analyse combinant données pluviométriques et imagerie satellitaire Landsat. Cette étude élargit la littérature récente de la seule mesure statistique vers l'étude des répercussions environnementales concrètes.

\enlargethispage{5\baselineskip}
En résumé, les revues théorique et empirique mettent en lumière la complexité des dynamiques climatiques et de leur variabilité pluviométrique au Sénégal. La revue théorique montre que cette variabilité trouve son explication dans l'articulation entre les forçages océaniques de grande échelle, l'équilibre radiatif global du système climatique et les outils statistiques qui permettent d'objectiver les ruptures observées dans les séries de données. Ces trois approches se complètent sans qu'aucune ne suffise isolément à rendre compte du phénomène.

La revue empirique fait ressortir un consensus autour de la rupture pluviométrique de la fin des années 1960, confirmée par des indicateurs aussi variés que les précipitations, les débits fluviaux et la dynamique du couvert forestier. Certains aspects demeurent toutefois discutés, notamment la nature exacte de la reprise pluviométrique amorcée dans les années 2000 et l'ampleur de la rupture selon les stations considérées. Les résultats apparaissent contrastés selon les zones étudiées et les méthodes mobilisées. Ces divergences soulignent l'importance de prendre en compte l'hétérogénéité spatiale et méthodologique dans l'analyse de la variabilité pluviométrique au Sénégal.